% =====================================================================================
% REVISIÓN 4  -  Conclusiones y Trabajo futuro  (la Revisión 4 exige TODAS las secciones anteriores terminadas)
% Escribe tu texto FUERA de los recuadros \begin{guia} ... \end{guia}.
% =====================================================================================

\section{Conclusiones} \label{sec:conclusiones} % audit:req=conclusiones

\begin{guia}[Guía: Conclusiones (Revisión 4)]
\textbf{¿Para qué sirve?} Cierra el trabajo: dice qué se logró, con qué evidencia y qué significa. Es lo primero que lee el
evaluador después del resumen.

\textbf{Debe responder:}
\begin{itemize}[nosep,leftmargin=*]
  \item ¿Se cumplió el objetivo general? ¿En qué medida y con qué evidencia (cita la sección, tabla o figura)?
  \item Para \emph{cada} objetivo específico: ¿se logró y dónde se demuestra?
  \item ¿Se resolvió el problema planteado? ¿Qué cambió para la empresa o los usuarios? Si hubo hipótesis: ¿se sustentó?
  \item ¿Cuáles fueron las principales dificultades y las lecciones aprendidas?
  \item ¿Qué limitaciones tienen los resultados (retoma las \emph{Limitaciones} de la Introducción)?
  \item ¿Qué aprendiste y cómo se relaciona con tu perfil de egreso?
\end{itemize}

\textbf{Extensión mínima:} 150 palabras.

\textbf{Reglas:} no introduzcas información nueva; no afirmes más de lo que los resultados sustentan; escribe en pasado
y con base en datos (``se redujo el tiempo de captura de 12 a 4 minutos'') y no en deseos.

\textbf{Errores frecuentes:} repetir el resumen; conclusiones genéricas (``se cumplió todo''); olvidar algún objetivo
específico; no mencionar limitaciones.

\textbf{Cómo se revisa:} correspondencia uno a uno con los objetivos, evidencia y honestidad.
\end{guia}

% >>> ESCRIBE AQUÍ las conclusiones <<<


\section{Trabajo futuro} \label{sec:trabajo_futuro} % audit:req=trabajo_futuro

\begin{guia}[Guía: Trabajo futuro (Revisión 4)]
\textbf{Debe responder:} ¿qué mejoras o extensiones concretas se podrían hacer? ¿qué quedó fuera del alcance y por qué?
¿qué recomiendas a la empresa o a quien continúe el trabajo? Prioriza (corto, mediano y largo plazo) y sé específico
(``agregar notificaciones por mensajería instantánea'', no ``mejorar el sistema'').

\textbf{Extensión mínima:} 50 palabras.
\end{guia}

% >>> ESCRIBE AQUÍ el trabajo futuro <<<
